% !TeX encoding = UTF-8
\documentclass[institutional,screen,none]{ua-engineering-v6-beamer}
% !TeX encoding = UTF-8
\UASetUniversity{Universidad Austral}
\UASetFaculty{Facultad de Ingeniería}
\UASetDepartment{Departamento de Computación}
\UASetCampus{Pilar}
\UASetCareer{Ingeniería Informática}
\UASetCommission{Comisión A}
\UASetProfessor{Equipo docente}
\UASetSemester{Segundo cuatrimestre 2026}
\UASetCourse{Sistemas Distribuidos}
\UASetDate{2026}


\UASetClassNumber{13}
\UASetClassTitle{Confiabilidad operativa}
\title{\UACourse}
\subtitle{Clase 13 --- SLI, SLO, SLA, error budgets y operación guiada por objetivos}

\author{\UAProfessor}
\date{\UADate}
\begin{document}

{
\setbeamertemplate{footline}{}
\begin{frame}[plain]
\titlepage
\end{frame}
}

\UASectionSlide{Del sistema al servicio}

\begin{frame}{La pregunta correcta}
Después de diseñar, observar y probar fallas, queda una pregunta:

\begin{block}{}
¿Qué significa, en términos operativos, que el sistema esté funcionando suficientemente bien?
\end{block}
\end{frame}

\begin{frame}{Confiabilidad no es perfección}
Un sistema confiable no es uno que:
\begin{itemize}
\item nunca falla
\item nunca degrada
\item siempre responde idealmente
\end{itemize}

\pause

Es uno que:
\begin{itemize}
\item cumple objetivos explícitos
\item dentro de una tolerancia aceptada
\end{itemize}
\end{frame}

\UASectionSlide{SLI}

\begin{frame}{SLI}
\begin{block}{}
Service Level Indicator = métrica que captura una propiedad relevante del servicio desde la perspectiva del usuario.
\end{block}

\pause

Ejemplos:
\begin{itemize}
\item porcentaje de requests exitosas
\item latencia p95
\item tiempo de procesamiento
\end{itemize}
\end{frame}

\begin{frame}{Un buen SLI}
Debe ser:
\begin{itemize}
\item medible
\item relevante para el usuario
\item estable
\item accionable
\end{itemize}
\end{frame}

\UASectionSlide{SLO}

\begin{frame}{SLO}
\begin{block}{}
Service Level Objective = objetivo cuantitativo sobre un SLI.
\end{block}

\pause

Ejemplo:
\begin{itemize}
\item 99.9\% de requests exitosas en 30 días
\item p95 de latencia menor a 300 ms
\end{itemize}
\end{frame}

\begin{frame}{Por qué importa}
Un SLO convierte una intuición vaga de “calidad” en:
\begin{itemize}
\item objetivo explícito
\item criterio de operación
\item base para priorización
\end{itemize}
\end{frame}

\UASectionSlide{SLA}

\begin{frame}{SLA}
\begin{block}{}
Service Level Agreement = compromiso formal con una parte externa, normalmente con implicancias contractuales.
\end{block}

\pause

Diferencia:
\begin{itemize}
\item SLI = métrica
\item SLO = objetivo interno
\item SLA = compromiso externo
\end{itemize}
\end{frame}

\UASectionSlide{Error Budget}

\begin{frame}{Error Budget}
\begin{block}{}
Si el SLO no exige perfección, entonces existe una cantidad tolerada de falla.
\end{block}

\pause

Ejemplo:
\begin{itemize}
\item SLO = 99.9\%
\item error budget = 0.1\%
\end{itemize}
\end{frame}

\begin{frame}{Interpretación}
El error budget no autoriza “hacer mal las cosas”.

\pause

Sirve para:
\begin{itemize}
\item equilibrar confiabilidad e innovación
\item decidir ritmo de cambios
\item guiar operación
\end{itemize}
\end{frame}

\UASectionSlide{Operación guiada por objetivos}

\begin{frame}{Sin SLOs}
Operar sin SLOs lleva a:
\begin{itemize}
\item discusiones vagas
\item prioridades arbitrarias
\item incidentes interpretados sin marco
\end{itemize}
\end{frame}

\begin{frame}{Con SLOs}
Operar con SLOs permite:
\begin{itemize}
\item medir desviación
\item decidir cuándo frenar cambios
\item justificar inversión en confiabilidad
\item alinear técnica y negocio
\end{itemize}
\end{frame}

\UASectionSlide{Diseño de buenos objetivos}

\begin{frame}{Errores típicos}
\begin{itemize}
\item elegir métricas internas irrelevantes para usuarios
\item fijar objetivos imposibles
\item confundir promedio con percentiles
\item mezclar múltiples propiedades en un solo indicador confuso
\end{itemize}
\end{frame}

\begin{frame}{Ejemplos de SLI/SLO}
\begin{center}
\begin{tabular}{ll}
\toprule
Servicio & SLO ejemplo \\
\midrule
Checkout & 99.95\% éxito en 30 días \\
API pública & p95 < 250 ms \\
Streaming & 99.9\% eventos procesados < 5 s \\
\bottomrule
\end{tabular}
\end{center}
\end{frame}

\UASectionSlide{Cierre}

\begin{frame}{Idea central}
\begin{block}{}
La confiabilidad operativa no se gestiona con intuiciones, sino con objetivos explícitos y presupuestos de error.
\end{block}
\end{frame}


\begin{frame}{Caso de diseño: decidir antes de nombrar tecnología}
\small
\begin{block}{Escenario}Un servicio crítico debe seguir operando ante latencia variable y fallas parciales. El equipo propone ``agregar más infraestructura'' sin declarar qué propiedad necesita preservar.\end{block}
\begin{enumerate}
\item ¿Cuál es el invariante o garantía que realmente importa?
\item ¿Qué falla concreta amenaza esa garantía?
\item ¿Qué mecanismo de Confiabilidad operativa es pertinente y cuál sería sobreingeniería?
\item ¿Qué métrica permitiría comprobar el costo de la decisión?
\end{enumerate}
\end{frame}
\begin{frame}{Checkpoint de comprensión}
\small
Al terminar esta clase deberías poder:
\begin{itemize}
\item explicar el problema sin apoyarte en nombres de productos;
\item identificar el supuesto de falla más importante;
\item distinguir una propiedad de seguridad de una de progreso;
\item construir un escenario donde la solución deja de ser suficiente;
\item defender el trade-off ante una objeción técnica.
\end{itemize}
\end{frame}

\begin{frame}{Burn rate}
No alcanza con saber cuánto budget queda: importa a qué velocidad se consume. Un burn rate alto exige reacción antes de agotar la ventana completa.
\end{frame}
\begin{frame}{Dos ventanas, una señal accionable}
Alertar con una ventana corta detecta incidentes rápidos; una ventana larga evita ruido. Combinar ambas reduce falsos positivos sin ignorar degradación sostenida.
\end{frame}
\begin{frame}{Checkpoint de servicio}
Defina un SLI desde la experiencia del usuario, su población válida, ventana, objetivo y política concreta cuando el budget se consume.
\end{frame}

\end{document}
